%% ma: L12-B13-0044
%% lop: 12
%% chuong: 4
%% bai: 13
%% dang: 12.13.23
%% loai: TL
%% muc: VD
%% dap_an: "-"
%% nguon: "(NBV)-10 ĐỀ KIỂM TRA HK2-TOÁN 12-DUNG CHUNG 3 SÁCH.pdf (H03, Đề số 3), Phần 4 Câu 1 (tự luận)"
%% kien_thuc: "Bài 13 (diện tích hình phẳng giới hạn bởi hai đồ thị), phương trình bậc ba có ba nghiệm cho trước (phân tích thành tích nhân tử, đồng nhất hệ số)"
%% trang_thai: chinh-thuc
%% ghi_chu: "Đáp án gốc 4 khớp. Hình dựng lại bằng hàm minh họa f(x) = x³/2 + 5x²/4 - x - 1/2 và g(x) = -x²/4 - x/2 + 1 (thỏa f - g = (x+3)(x²-1)/2; chỉ hiệu f - g được xác định bởi đề, hình không theo tỉ lệ ảnh gốc). Đã duyệt 21/10/2026 theo yêu cầu của giáo viên: giải lại độc lập, kiểm bằng sympy, soát trùng, biên dịch OK."
\begin{ex}
Cho hai hàm số $f(x)=ax^3+bx^2+cx-\dfrac12$ và $g(x)=dx^2+ex+1$ $(a,b,c,d,e\in\mathbb R)$. Biết rằng đồ thị của hàm số $y=f(x)$ và $y=g(x)$ cắt nhau tại $3$ điểm có hoành độ lần lượt là $-3;-1;1$ (tham khảo hình vẽ). Tính diện tích hình phẳng giới hạn bởi hai đồ thị đã cho.
\begin{center}
\begin{tikzpicture}[xscale=1.1,yscale=.9]
  \fill[gray!45] plot[domain=-3:-1,samples=40] (\x,{.5*\x*\x*\x+1.25*\x*\x-\x-.5}) -- plot[domain=-1:-3,samples=40] (\x,{-.25*\x*\x-.5*\x+1}) -- cycle;
  \fill[gray!45] plot[domain=-1:1,samples=40] (\x,{.5*\x*\x*\x+1.25*\x*\x-\x-.5}) -- plot[domain=1:-1,samples=40] (\x,{-.25*\x*\x-.5*\x+1}) -- cycle;
  \draw[thick,domain=-3.4:1.35,samples=70] plot (\x,{.5*\x*\x*\x+1.25*\x*\x-\x-.5});
  \draw[thick,domain=-3.6:1.9,samples=70] plot (\x,{-.25*\x*\x-.5*\x+1});
  \draw[truc] (-3.9,0)--(2.2,0) node[below]{$x$};
  \draw[truc] (0,-1.9)--(0,3.1) node[left]{$y$};
  \draw[phu] (-3,0)--(-3,.25) (-1,0)--(-1,1.25) (1,0)--(1,.25);
  \path (0,0) node[below right]{\footnotesize$O$} (-3,0) node[below right]{\footnotesize$-3$} (-1,0) node[below]{\footnotesize$-1$} (1,0) node[below]{\footnotesize$1$};
  \path (-2.0,2.55) node[above left]{\footnotesize$y=f(x)$} (1.4,0.4) node[right]{\footnotesize$y=g(x)$};
\end{tikzpicture}
\end{center}
\loigiai{Phương trình hoành độ giao điểm: $f(x)-g(x)=ax^3+(b-d)x^2+(c-e)x-\dfrac32=0$ có ba nghiệm $-3;-1;1$ nên $f(x)-g(x)=a(x+3)(x+1)(x-1)$.

So sánh hệ số tự do: $-3a=-\dfrac32\Rightarrow a=\dfrac12$. Vậy $f(x)-g(x)=\dfrac12\left(x^3+3x^2-x-3\right)$.

Diện tích: $S=\displaystyle\int_{-3}^{1}|f(x)-g(x)|\,\mathrm dx$. Trên $[-3;-1]$ có $f(x)\ge g(x)$, trên $[-1;1]$ có $f(x)\le g(x)$ nên
\[S=\int_{-3}^{-1}\dfrac12\left(x^3+3x^2-x-3\right)\mathrm dx-\int_{-1}^{1}\dfrac12\left(x^3+3x^2-x-3\right)\mathrm dx=2+2=4.\]}
\end{ex}
